\section{Threats to validity}
\label{sec:limits}

\subsection{Construct validity}

\paragraph{Our own definition of scope.} The primary measure counts edits to
four functions that we chose, and the broader measure adds tests, other
functions, and new files, but neither catches every possible excursion (for
example, a changed docstring in the target function). A reader could argue that, under the loose
prompt, some of those edits were invited. We think the more useful reading is
the one a reviewer of a bug-fix pull request would take, but we report the
hidden-test outcome alongside the symbol count so the two readings can be
separated.

\paragraph{Arms, not models.} Every comparison is between shipped
configurations. Within the Cursor arms the scaffold is shared, which narrows
the gap between an arm and its model, but the scaffold may still interact
with each model differently, and the single Claude Code arm in generation~1
differs from the rest in both model and harness.

\paragraph{What was recorded about each model.} We record the model flag we
passed to each client, not the model version the vendor actually served or
the client version, so we cannot rule out silent changes on the vendor side
during the study.

\subsection{Internal validity}

\paragraph{Agents were not sandboxed.} Each agent started in a fresh copy of
the task, but that copy sat inside the benchmark directory and the agents
had ordinary file-system access. Some agents walked up the tree and read
canonical task files, hidden tests, reference solutions, or earlier runs. We
detect these runs from their trajectories and exclude them, and the main
results are the same with or without them, but the detection rule can miss
access that leaves no trace in the tool calls. Future runs should place the
working copy outside the benchmark tree or in a container.

\paragraph{The working directory named the rung.} Each run's directory name included the task slug, such as
\texttt{T005-percentile-loose-prompt}, and an agent could see it. We have no
evidence that agents used this, but it is a channel through which an agent
could learn which condition it was in.

\paragraph{Rungs differ in more than register.} The scoped prompts end with
the command that runs the visible tests; only \NTestCmdLoose{} of the
\NDomains{} loose prompts and \NTestCmdVague{} vague prompts do. This could raise
exploration on the loose and vague rungs, so we do not read the differences in
unique tool calls as a pure effect of wording. It is harder to see how it
could explain planted-function edits under the loose prompt alone, since the
vague prompts lack the command as well.

\paragraph{Incomplete and long sessions.} We treat a run as valid if the
harness recorded no error and the agent made at least one tool call. In the
analyzed set, \XIncomplete{} valid trajectories end without a final result
event and \XOverTime{} runs recorded a wall time above the nominal 600-second limit
without being flagged as timeouts, for reasons we have not diagnosed. We kept these runs, since they produced a scorable working
copy, but a stricter validity rule would drop them.

\paragraph{Lost runs.} Some generation-2 arms lost runs to vendor usage caps
or to sessions that ended before the agent acted. These invalid runs are
excluded rather than scored, and the counts are reported per arm
(Table~\ref{tab:coverage}). If losses were related
to task content, the remaining sample could be biased; we saw no sign of
this. Nearly all invalid runs were started during a few hours on a single
evening (28 September 2026), and retries of the same tasks in October
completed normally, which points to a temporary vendor-side condition rather
than to anything about the tasks.

\paragraph{Few runs per cell.} Most arm-task cells hold one run. Agent
behavior is not deterministic, so an individual cell (one arm on one loose
task) should be read as one draw rather than as that arm's behavior. Our
conclusions rest on pooled rates over many cells and on the within-task
contrast between rungs, which is less exposed to this noise than any single
cell.

\subsection{External validity}

\paragraph{Narrow tasks.} The repositories are small, the bugs are local, and
every task is in Python. This is deliberate, because it makes ``out of scope''
well defined, but it means our rates describe a clean laboratory setting.
Real repositories offer far more to drift into, so the absolute rates there
may differ; we have not measured them.

\paragraph{Prompt calibration from one person.} The vague rung was written with
the prompting habits of one developer in mind, and the design-time profile we
report comes from that developer's private history. Other developers may write longer or
more specific requests, and the loose and vague templates are our own
wording. A broader sample of real prompts would make the ladder more
representative.
